\emph{证明系统}提供刻画语句间后承关系及相容性的纯句法方法。
\emph{自然演绎}就是这样一种证明系统。
其中的\emph{推导}是一棵由公式构成的树。
推导中最顶端的公式是\emph{假设}。
为使推导正确，所有其他公式都必须由若干\emph{推理规则}中的某一条正确证成。
这些规则成对出现：每个联结词和量词都有一条引入规则和一条消去规则。
例如，若公式~$!A$ 由 $\Elim{\lif}$（蕴涵消去）规则证成，
则它前面的公式（即\emph{前提}）必须是 $!B \lif !A$ 和 $!B$（对某个~$!B$）。
有些推理规则还允许\emph{解除}假设。
例如，若使用 $\Intro{\lif}$（蕴涵引入）规则从 $!B$ 推出 $!A \lif !B$，
则在通向前提~$!B$ 的推导中，所有作为假设出现的~$!A$ 均可被解除，
并被赋予一个标签，该标签同时记录在该推理步骤处。

如果存在一个以公式~$!A$ 为末端公式且所有假设均被解除的推导，
我们就称 $!A$ 是一个\emph{定理}，并记作~$\Proves !A$。
如果所有未被解除的假设都属于某个集合~$\Gamma$，
我们就称 $!A$ \emph{可从}~$\Gamma$ \emph{推导}，并记作 $\Gamma \Proves !A$。
如果 $\Gamma \Proves \lfalse$，则称 $\Gamma$ 是\emph{不一致的}，否则称为\emph{一致的}。
这些概念相互关联，例如，$\Gamma \Proves !A$ 当且仅当 $\Gamma \cup \{\lnot !A\}$ 不一致。
它们也与相应的语义概念相关，例如，若 $\Gamma \Proves !A$，则 $\Gamma \Entails !A$。
证明系统的这一性质——即凡从 $\Gamma$ 可推导出的都保证是 $\Gamma$ 的语义后承——称为\emph{可靠性}。
\emph{可靠性定理}通过对推导长度进行归纳来证明，其要点在于表明：每个推理步骤都能保持如下性质——若前提为各自未解除假设的语义后承，则结论亦为开放假设的语义后承。